\section{Cierre y fuentes}

\begin{frame}{Ideas que deben quedar}
\footnotesize
\begin{columns}[T,onlytextwidth]
\begin{column}{0.60\textwidth}
\begin{enumerate}
  \item Una SVM lineal busca el hiperplano con máximo margen.
  \item Hard margin exige separación perfecta; soft margin introduce holguras y regularización.
  \item En scikit-learn, $C$ alto penaliza más los errores; $C$ bajo regulariza más.
  \item La formulación dual depende de productos internos, lo que habilita kernels.
  \item El kernel trick equivale a trabajar en un espacio de características ampliado sin construirlo explícitamente.
  \item En RBF, $\gamma$ gobierna la localidad; $C$ y $\gamma$ deben validarse conjuntamente, siempre dentro de un pipeline con escalamiento y validación cruzada.
\end{enumerate}
\end{column}
\begin{column}{0.36\textwidth}
\begin{keybox}
\centering
\large
\textbf{Margen}\\[0.14cm]
$+$\\[0.14cm]
\textbf{Regularización}\\[0.14cm]
$+$\\[0.14cm]
\textbf{Kernel}\\[0.18cm]
$\Downarrow$\\[0.18cm]
\textbf{Fronteras flexibles con optimización convexa}
\end{keybox}
\end{column}
\end{columns}
\end{frame}

% ================================================================
\begin{frame}{Bibliografía base del Drive}
\footnotesize
\begin{ccdebox}
\textbf{Géron, A.} (2022). \emph{Hands-On Machine Learning with Scikit-Learn, Keras, and TensorFlow}, 3.ª ed., O'Reilly.\\
\textbf{Capítulo 5:} clasificación SVM lineal, soft margin, kernels, RBF, complejidad y formulación matemática.
\end{ccdebox}
\vspace{0.07cm}
\begin{ccdebox}
\textbf{James, G.; Witten, D.; Hastie, T.; Tibshirani, R.} (2013). \emph{An Introduction to Statistical Learning with Applications in R}. Springer.\\
\textbf{Capítulo 9:} maximal margin classifier, support vector classifier, kernels, RBF, multiclase y laboratorio.
\end{ccdebox}
\vspace{0.07cm}
\begin{ccdebox}
\textbf{Hastie, T.; Tibshirani, R.; Friedman, J.} (2009). \emph{The Elements of Statistical Learning}, 2.ª ed., Springer.\\
\textbf{Capítulo 12:} formulación avanzada de SVM, penalización, reproducing kernels y perspectiva estadística.
\end{ccdebox}
\end{frame}

% ================================================================
\begin{frame}{Material reproducible y compilación}
\small
\begin{columns}[T,onlytextwidth]
\begin{column}{0.48\textwidth}
\begin{ccdebox}
\textbf{Código asociado}\\[0.10cm]
\texttt{codigo\_demo\_svm.py}\\[0.10cm]
Genera las métricas y las figuras PNG utilizadas por la presentación.
\end{ccdebox}
\vspace{0.14cm}
\begin{keybox}
\texttt{pip install numpy matplotlib scikit-learn}\\[0.08cm]
\texttt{python codigo\_demo\_svm.py}
\end{keybox}
\end{column}
\begin{column}{0.48\textwidth}
\begin{ccdebox}
\textbf{Presentación Beamer}\\[0.10cm]
\texttt{presentacion-svm.tex}
\end{ccdebox}
\vspace{0.14cm}
\begin{keybox}
\texttt{pdflatex presentacion-svm.tex}\\[0.08cm]
\texttt{pdflatex presentacion-svm.tex}
\end{keybox}
\vspace{0.14cm}
\begin{ccdebox}
\footnotesize
El logo y el diseño provienen de la plantilla canónica \texttt{assets/ccde-beamer.sty}; si el logo no está disponible, el documento compila con un fallback textual.
\end{ccdebox}
\end{column}
\end{columns}
\end{frame}
